\section{Benchmark Construction and Experimental Protocol}
\label{app:data}
\label{app:protocol}
\subsection{Contract structures and paired boundaries}
The primary dataset contains 20 templates (Table~\ref{tab:templates}), with five feasible/infeasible pairs per template. Each purchase has six components; the retained events are the 57 affected subsets of size at most four, including the empty event. Empty events have zero assessment. A nonempty event incurs one administrative fee, outside insurance or compensation offsets. Public contracts state that no additional hidden clauses apply.

The master construction seed is 39071007, with predetermined template and stratum offsets. The unique authorized positive-value candidate has loss margin $m$: the authorized cap minus its worst loss, in subunits (Table~\ref{tab:strata}). Increasing only its administrative fee by $m+1$ makes the paired task infeasible, with maximum loss of 200,001 subunits against a 200,000-subunit limit. The resulting violation is always one subunit, although the fee increment depends on the feasible margin. Other candidates already exceed the limit by 1--20,000 subunits. Opaque candidate identifiers contain no answer labels, and positions are sampled before attributes are generated. Amounts are reported in MU without changing the stored values.

\begin{table*}[!tbp]
\centering\small
\begin{tabular}{llp{10.4cm}}
\toprule
ID & Family & Distinct contractual structure \\
\midrule
a1 & Coupon cancellation & Overlapping three-person groups; a clawback applies to partial but not full group cancellation.\\
a2 & Coupon cancellation & Different settlement charges for one versus both members of a two-person group.\\
a3 & Coupon cancellation & A cap on group clawbacks plus a conditional order-level clawback.\\
a4 & Coupon cancellation & A rebate offsets group charges only after multiple groups are broken.\\
a5 & Coupon cancellation & Only the largest group clawback is charged, with a specified combination surcharge.\\
\midrule
b1 & Insured cancellation & Item deductibles and item payout limits, followed by an aggregate payout cap.\\
b2 & Insured cancellation & An order-level deductible followed by proportional coinsurance.\\
b3 & Insured cancellation & The number of cancellations changes the aggregate insurance limit.\\
b4 & Insured cancellation & Primary insurance followed by secondary insurance on remaining loss.\\
b5 & Insured cancellation & Payout is triggered by a strict threshold, then calculated from the first unit of loss.\\
\midrule
c1 & Threshold returns & Two product groups have separate retained-value discount clawbacks.\\
c2 & Threshold returns & Fixed shipping parcels; return freight is the largest item fee within each affected parcel.\\
c3 & Threshold returns & Retained purchase value determines a tiered retained discount.\\
c4 & Threshold returns & Separate retained-value and retained-weight thresholds govern two subsidies.\\
c5 & Threshold returns & Restocking fees are capped; sufficiently large returns waive return freight.\\
\midrule
d1 & Service replacement & Item-level compensation limits apply before an order-level cap.\\
d2 & Service replacement & Activating replacement service incurs a fixed fee per affected group.\\
d3 & Service replacement & The number of failures determines the aggregate compensation cap.\\
d4 & Service replacement & Cash compensation combines with explicitly authorized, discounted service credits.\\
d5 & Service replacement & Overlapping group dispatch fees and threshold-triggered additional compensation.\\
\bottomrule
\end{tabular}
\caption{All 20 primary structures. Numerical instances vary within each template; the four families contain 50 cases each. These contracts define synthetic assessment rules and are not claims about actual merchants' terms.}
\label{tab:templates}
\end{table*}

\begin{table}[!htbp]
\centering\small
\begin{tabular}{rrrr}
\toprule
Stratum & Candidates & Margin (subunits) & Cases\\
\midrule
1 & 128 & 1 & 40\\
2 & 256 & 100 & 40\\
3 & 512 & 10,000 & 40\\
4 & 1,024 & 100 & 40\\
5 & 1,024 & 1 & 40\\
\bottomrule
\end{tabular}
\caption{Each stratum includes one pair for every template. Catalogue size, margin, and seed change together; these strata do not isolate their individual causal effects.}
\label{tab:strata}
\end{table}

The design fixes unique feasibility, balanced acceptance/refusal, and complete catalogues. Size, margin, and seed vary together across strata, so their individual effects cannot be estimated separately. The task balance is not an estimate of real purchasing frequencies.

\subsection{Reference calculations and freezing}
A cash-flow ledger evaluates bitmask events with integer or rational arithmetic, and a separate expression implementation uses the rule engine. Before model evaluation, we compare candidate feasibility, worst loss, guaranteed utility, and final decisions. They agree on all 6,712,320 candidate--event checks. Twenty fresh attribute assignments per template are also checked on all 64 subsets, yielding $20\times20\times64=25{,}600$ checks, plus three hand-calculated threshold witnesses. Both implementations arose within one development workflow and may share a misunderstanding of a clause despite agreeing numerically.

Tasks, references, prompts, and evaluator code are frozen before model evaluation. Public inputs exclude construction seeds, reference matrices, labels, and designated safe identifiers. Candidate visibility follows Table~\ref{tab:visibility}.

\subsection{Common instructions and outcome protocol}
\label{app:implementation}
All primary routes receive the same substantive policy, summarized below in English. The frozen model-facing prompts retain their original language and are unchanged by this summary.
\begin{quote}\small
Retain applicable interpretations. Respect current authorization in every retained state. Evaluate hard requirements and the complete risk permission, then maximize the minimum state-dependent utility. Buy only for positive guaranteed value; otherwise refuse. Do not ask the user or change permissions.
\end{quote}
Direct returns a decision, allowed identifiers, and guaranteed utility. Set returns the state-wise expressions of Appendix~\ref{app:interface}. Exactness requires agreement with the reference decision and, for acceptance, a nonempty subset of authorized maximizers and correct value. Disjoint failures are unsafe acceptance, false refusal, and Other (remaining value or invalid-execution failures). Set executions must satisfy the complete-domain checks in Algorithm~\ref{alg:decision}. Failed or unresolved execution counts as Other, including a no-action output whose reference decision is refusal.

\begin{table}[!htbp]
\centering\small
\begin{tabular}{lcc}
\toprule
Property & Direct & Set\\
\midrule
Natural-language records & Yes & Yes\\
Attribute schema & Yes & Yes\\
Catalogue in model input & Yes & No\\
Deterministic row evaluation & No & Yes\\
Exact arithmetic in executor & --- & Yes\\
Observable allowed domain & No & Yes\\
Model calls per primary trial & 1 & 1\\
\bottomrule
\end{tabular}
\caption{Primary method controls. Different catalogue visibility is a declared architectural choice. ``One call'' refers to a logical trial; transport retries are separately retained.}
\label{tab:visibility}
\end{table}

A primary comparison trial allows one semantic attempt. Transport failures may be retried within the recorded limit; formatting, semantic, and output-budget failures end the trial. We save every received response before scoring and count all failures in completed cohorts. Calls not issued remain missing. We reuse an earlier result only when the messages, schema, requested model, effort setting, deadline, output budget, and both task and reference data are unchanged. Sources are selected without using the outcome. The availability study's first-evaluable retention rule is specified in Appendix~\ref{app:availability}.

\subsection{Model configurations}
\label{app:model-configurations}
\label{app:qwen}
\label{app:primary-extensions}
Table~\ref{tab:configurations} specifies the primary and intent settings. Its continuation lists 49 model--setting combinations across 12 requested identities. Exploratory settings document configuration coverage and do not add to the primary trial count. Paired Direct/Set comparisons share requested generation settings. Primary command-line interface (CLI) calls use fresh contexts with tools and delegation disabled; application programming interface (API) outputs pass the common validator. Requested effort and output caps do not establish equivalent compute across providers.

\begin{table*}[!tbp]
\centering\small
\setlength{\tabcolsep}{4pt}
\begin{tabular}{@{}llrr@{}}
\toprule
Configuration & Requested model & Output cap & Deadline\\
\midrule
\multicolumn{4}{@{}l}{\textbf{Primary 200-task comparison}}\\
DeepSeek low / high & \texttt{deepseek-flash} & 262,144 & 1,800\\
Qwen Flash low & \texttt{qwen3.8-flash} & 131,072 & 1,800\\
Qwen Max low & \texttt{qwen3.8-max} & 131,072 & 1,800\\
GLM Flash low & \texttt{glm-5.3-flash} & 131,072 & 1,800\\
GLM high & \texttt{glm-5.3} & 131,072 & 1,800\\
5.6 Luna medium & \texttt{gpt-5.6-luna} & 128,000 & 1,800\\
5.6 Terra medium & \texttt{gpt-5.6-terra} & 128,000 & 1,800\\
5.6 Sol medium & \texttt{gpt-5.6-sol} & 128,000 & 1,800\\
Astra max & \texttt{gpt-6-astra} & 128,000 & 1,800\\
Opus 5.5 medium & \texttt{claude-opus-5-5} & 128,000 & 1,800\\
\midrule
\multicolumn{4}{@{}l}{\textbf{Retained-intent comparison (Figure~\ref{fig:intent})}}\\
DeepSeek low / max & \texttt{deepseek-flash} & 32,768 & 300\\
GLM low & \texttt{glm-5.3} & 32,768 & 300\\
Qwen Max low & \texttt{qwen3.8-max-0902} & 32,768 & 300\\
\midrule
\multicolumn{2}{@{}l}{\textbf{Prospective availability (Figure~\ref{fig:ablations})}} & \multicolumn{2}{r@{}}{Output cap}\\
DeepSeek low & \texttt{deepseek-flash} & \multicolumn{2}{r@{}}{32,768}\\
GLM low & \texttt{glm-5.3} & \multicolumn{2}{r@{}}{32,768}\\
Qwen Max low & \texttt{qwen3.8-max} & \multicolumn{2}{r@{}}{32,768}\\
\bottomrule
\end{tabular}

\caption{Requested model configurations. Output caps are in tokens and deadlines in seconds. Each named primary configuration uses 200 tasks per method. The lower block gives the four configurations in Figure~\ref{fig:intent}; their task counts appear in Table~\ref{tab:intent-figure3}. Configuration names refer to requests, not independently attested backend versions.}
\label{tab:configurations}
\end{table*}

\begin{table*}[!tbp]
\addtocounter{table}{-1}
\centering\small
\setlength{\tabcolsep}{4pt}
\begin{tabular}{@{}>{\raggedright\arraybackslash}p{.30\textwidth}>{\raggedright\arraybackslash}p{.255\textwidth}>{\raggedright\arraybackslash}p{.23\textwidth}>{\raggedright\arraybackslash}p{.135\textwidth}@{}}
\toprule
Configuration & Requested model & Output caps & Deadlines\\
\midrule
\textbf{DeepSeek Flash}\newline off, low, high, max & \texttt{deepseek-flash} & 8,192, 16,384, 24,000, 32,768, 48,000, 65,536, 262,144 & 180, 240, 300, 360, 420, 600, 900, 1,800\\
\addlinespace[5pt]
\textbf{DeepSeek V4 Pro}\newline off, low, high, max & \texttt{deepseek-v4-pro} & 32,768 & 300\\
\addlinespace[5pt]
\textbf{Qwen Flash}\newline off, low, medium, xhigh & \texttt{qwen3.8-flash} & 24,000, 32,768, 131,072 & 300, 420, 1,800\\
\addlinespace[5pt]
\textbf{Qwen Max}\newline low, xhigh & \texttt{qwen3.8-max} & 131,072 & 1,800\\
\addlinespace[5pt]
\textbf{Qwen Max (0902)}\newline off, low, medium, xhigh & \texttt{qwen3.8-max-0902} & 32,768 & 300\\
\addlinespace[5pt]
\textbf{GLM Flash}\newline low, high, max & \texttt{glm-5.3-flash} & 24,000, 32,768, 65,536, 131,072 & 300, 420, 1,800\\
\addlinespace[5pt]
\textbf{GLM}\newline low, high, max & \texttt{glm-5.3} & 32,768, 65,536, 131,072 & 300, 1,800\\
\addlinespace[5pt]
\textbf{5.6 Luna}\newline off, low, medium, high, xhigh, max & \texttt{gpt-5.6-luna} & 128,000 & 1,800\\
\addlinespace[5pt]
\textbf{5.6 Terra}\newline off, low, medium, high, xhigh, max & \texttt{gpt-5.6-terra} & 128,000 & 1,800\\
\addlinespace[5pt]
\textbf{5.6 Sol}\newline off, low, medium, high, xhigh, max & \texttt{gpt-5.6-sol} & 128,000 & 1,800\\
\addlinespace[5pt]
\textbf{Astra}\newline low, medium, high, xhigh, max & \texttt{gpt-6-astra} & 128,000 & 300, 600, 900, 1,800\\
\addlinespace[5pt]
\textbf{Opus 5.5}\newline medium, max & \texttt{claude-opus-5-5} & 128,000 & 900, 1,800\\
\addlinespace[5pt]
\bottomrule
\end{tabular}

\caption{Requested model configurations (continued): recorded model--setting coverage. Off denotes the native \texttt{none} setting. Output caps are in tokens and deadlines in seconds. Multiple values summarize settings observed in separate studies; they are not a fully crossed experimental grid. The primary and retained-intent settings are specified in the preceding part. Performance comparisons retain their own study cohorts.}
\label{tab:configurations-extended}
\end{table*}

\subsection{Unified reproduction entry point}
\label{app:reproduction}
The anonymous artifact provides implementation code, benchmark inputs and references, prompts, generation settings, retained outputs, and reporting materials at \url{https://anonymous.4open.science/r/risk-aware-payment-delegation-D4FC/}. Its documentation gives the environment setup and commands for rebuilding the reported results and running the corresponding checks. Table~\ref{tab:artifact-map} distinguishes report reconstruction, retained-output verification, and reference recomputation. These offline checks make no model calls and leave experimental records unchanged; new inference runs are documented separately.

\begin{table*}[t]
\centering\small\setlength{\tabcolsep}{4pt}
\begin{tabular}{@{}p{.20\textwidth}p{.39\textwidth}p{.36\textwidth}@{}}
\toprule
Result & Supporting materials & Verification level\\
\midrule
Primary outcomes & Benchmark instances, per-trial observations, and separately indexed aggregate updates & Rebuild tables from their designated inputs; reference checking is separate.\\
Retained requests and decisions & Original requests, responses, configurations, and stored evaluations & Check response integrity and agreement with stored scores; rescore retained prompt-control decisions.\\
Intent and observed trace & Paired tasks, original generated conditions, and the evidence map & Rebuild intent counts and verify the displayed traces against retained outputs.\\
Replay and interventions & Frozen specifications, catalogue conditions, and recorded interventions & Recompute diagnostic aggregates with collection settings and per-output provenance retained.\\
Decision contract and plots & Evaluator, reference ledger, reporting inputs, and figure-generation tools & Check Algorithm~\ref{alg:decision} and regenerate figures and tables; reference recomputation is a separate check.\\
\bottomrule
\end{tabular}
\caption{\textbf{Result-to-artifact map.} The artifact documentation links each result to its supporting materials and verification workflow. Report reconstruction, response rescoring, and new model calls are separate procedures. Aggregate updates remain separate from archived per-trial observations.}
\label{tab:artifact-map}
\end{table*}
